\documentclass[11pt,a4paper]{article}

% --- Page margins and paragraphs ---
\usepackage[margin=1in]{geometry}
\usepackage{parskip}

% --- Basic math and typesetting ---
\usepackage{amsmath,amssymb,amsthm}
\usepackage{xcolor} % 

% --- Theorem and definition environments ---
\theoremstyle{definition}
\newtheorem{definition}{Definition}[section]
\newtheorem{example}{Example}[section]

\theoremstyle{plain}
\newtheorem{theorem}{Theorem}[section]
\newtheorem{proposition}{Proposition}[section]
\newtheorem{lemma}{Lemma}[section]

\begin{document}

% ==================================================
% Title and author information
% ==================================================
\begin{center}
    {\LARGE \textbf{Mathematics Method for Physics--Lecture 16}} \\[0.6em]
    {\large Laplace Transform: Definition and Properties}
\end{center}

\vspace{-0.5em}
\noindent\rule{\textwidth}{0.5pt}
\vspace{0.3em}

\begin{center}
    Instructor: Zhaolong Han, SOCE, Shanghai Jiao Tong University;
\end{center}

\vspace{-0.6em}
\noindent Email: han.arkey@sjtu.edu.cn

\vspace{0.4em}
\noindent\rule{\textwidth}{0.5pt}

\vspace{0.8em}
\noindent\textit{Conventions: $p=\sigma+\mathrm{i}\omega$; transform pairs are written function
first, $\{f(t),F(p)\}$; $\mathrm{e},\mathrm{i},\mathrm{d}$ upright; formulas are numbered
(16.1)--(16.24).}

\vspace{1.5em}

\setcounter{section}{15}

% ==================================================
% Main content
% ==================================================
\section{Laplace Transform}

\subsection{Definition}

For $f(t)$ defined on $[0,+\infty)$ put
\begin{equation}
F(p)=\int_{0}^{+\infty} f(t)\,\mathrm{e}^{-pt}\,\mathrm{d}t, \tag{16.1}
\end{equation}
where
\begin{equation}
p=\sigma+\mathrm{i}\omega \tag{16.2}
\end{equation}
is a \textbf{complex variable}. Then $\{f(t),F(p)\}$ is a \textbf{Laplace transform pair}, written
\begin{equation}
\mathcal{L}\left[f(t)\right]=F(p),\qquad\text{or}\qquad
\mathcal{L}^{-1}\left[F(p)\right]=f(t). \tag{16.3}
\end{equation}

\textbf{Remarks}
\begin{enumerate}
    \item $F(p)$ is complex-valued. The imaginary part $\omega$ is exactly the Fourier frequency,
          the real part $\sigma$ acts as a \textbf{damping factor}: multiplying $f$ by
          $\mathrm{e}^{-\sigma t}$ before transforming is why the Laplace transform reaches
          functions of moderate growth that the Fourier transform does not.
    \item The lower limit is written $0$; for a function with an impulse at $t=0$ the left limit
          $0^{-}$ is meant --- with $0^{-}$ one gets $\mathcal{L}[\delta]=1$, with $0^{+}$ one
          gets $0$.
    \item Only $t\ge 0$ enters (16.1), so the values of $f(t)$ for $t<0$ do not affect $F(p)$:
          what is transformed is the causal part $f(t)u(t)$.
\end{enumerate}

\clearpage
\subsection{Exponential order and the existence theorem}

\subsubsection{Functions of exponential order}

\begin{definition}[Exponential order]
If there exist constants $M>0$ and $\sigma_c\ge 0$ such that
\begin{equation}
\left|f(t)\right|\le M\,\mathrm{e}^{\sigma_c t},\qquad t\ge 0, \tag{16.4}
\end{equation}
then $f(t)$ is said to be of \textbf{exponential order}, and $\sigma_c$ is called its
\textbf{growth exponent}.
\end{definition}

\begin{itemize}
    \item $u(t)$, $t^n$, $\mathrm{e}^{at}$, $\sin\beta t$, $\cos\beta t$ are all of exponential
          order;
    \item $t^{-1/2}$ is unbounded as $t\to 0^{+}$ and is \textbf{not} of exponential order, yet
          $\int_{0}^{+\infty} t^{-1/2}\mathrm{e}^{-pt}\,\mathrm{d}t$ still converges (integrable
          singularity);
    \item $\mathrm{e}^{t^{2}}$ grows faster than any exponential function and is not of
          exponential order --- its Laplace transform does not exist. Exponential order is
          therefore \textbf{sufficient, not necessary}.
\end{itemize}

\subsubsection{Theorem 16.1 (existence theorem)}

\begin{theorem}
Let $f(t)$ on $[0,+\infty)$ satisfy
\begin{enumerate}
    \item the \textbf{Dirichlet conditions} (piecewise continuous on every finite interval,
          finitely many discontinuities of the first kind, finitely many extrema), and
    \item $\left|f(t)\right|\le M\,\mathrm{e}^{\sigma_c t}$ (exponential order).
\end{enumerate}
Then
\begin{equation}
F(p)=\int_{0}^{+\infty} f(t)\,\mathrm{e}^{-pt}\,\mathrm{d}t \tag{16.5}
\end{equation}
exists and is \textbf{analytic} on the half-plane $\operatorname{Re}p>\sigma_c$.
\end{theorem}

\subsubsection{Proof}

\textbf{1. Absolute convergence.} For $\sigma=\operatorname{Re}p>\sigma_c$, using
$\left|\mathrm{e}^{-pt}\right|=\mathrm{e}^{-\sigma t}$,
\begin{equation}
\int_{0}^{+\infty}\left|f(t)\,\mathrm{e}^{-pt}\right|\mathrm{d}t
\le \int_{0}^{+\infty} M\,\mathrm{e}^{(\sigma_c-\sigma)t}\,\mathrm{d}t
=\frac{M}{\sigma-\sigma_c}<+\infty . \tag{16.6}
\end{equation}

\textbf{2. Uniform convergence and differentiability.} On $\operatorname{Re}p\ge\sigma_0>\sigma_c$
the integral (16.5) converges uniformly, and so does the differentiated integral
$\int_{0}^{+\infty} f(t)\,(-t)\,\mathrm{e}^{-pt}\,\mathrm{d}t$; differentiation under the integral
sign is therefore legitimate:
\begin{equation}
F'(p)=\int_{0}^{+\infty} f(t)\left(\mathrm{e}^{-pt}\right)'_{p}\mathrm{d}t
=-\int_{0}^{+\infty} t\,f(t)\,\mathrm{e}^{-pt}\,\mathrm{d}t . \tag{16.7}
\end{equation}

So $F$ is differentiable --- hence analytic --- on $\operatorname{Re}p>\sigma_c$.
\hfill$\blacksquare$

\subsubsection{Remarks}

\begin{enumerate}
    \item $F(p)$ is analytic for $\operatorname{Re}p>\sigma_c$, and
\begin{equation}
\lim_{\operatorname{Re}p\to+\infty}F(p)=0 . \tag{16.8}
\end{equation}
    \item By \textbf{analytic continuation} $F(p)$ is analytic on the whole $p$-plane except at
          its singularities; the positions of the poles / branch points determine the structure
          of $f(t)$.
    \item \textbf{Region of convergence.} The theorem gives $\operatorname{Re}p>\sigma_c$, i.e.\ the right half-plane $\sigma>\sigma_c$
          ($\sigma$ horizontal, $\mathrm{i}\omega$ vertical in the $p$-plane). It depends on
          $f$: $\mathcal{L}[u]$ converges for $\operatorname{Re}p>0$, $\mathcal{L}[\mathrm{e}^{at}]$ for
          $\operatorname{Re}p>\operatorname{Re}a$; for a rational image the ROC is fixed by the
          rightmost singularity.
\end{enumerate}

\clearpage
\subsection{Common transform pairs}

Every entry below is obtained by \textbf{direct integration} of (16.1).

\subsubsection{$u(t)\leftrightarrow 1/p$}

With $u(t)=1$ for $t\ge 0$ and $0$ for $t<0$,
\begin{equation}
\mathcal{L}\left[u(t)\right]=\int_{0}^{+\infty}\mathrm{e}^{-pt}\,\mathrm{d}t
=\left.\frac{\mathrm{e}^{-pt}}{-p}\right|_{0}^{+\infty}=\frac{1}{p},
\qquad \operatorname{Re}p>0 . \tag{16.9}
\end{equation}

\subsubsection{Powers $t$, $t^{2}$, \dots, $t^{n}$}

By parts ($u=t$, $\mathrm{d}v=\mathrm{e}^{-pt}\mathrm{d}t$):
\begin{equation}
\begin{split}
\mathcal{L}\left[t\right]&=\int_{0}^{+\infty} t\,\mathrm{e}^{-pt}\,\mathrm{d}t
=\left[t\left(-\frac{\mathrm{e}^{-pt}}{p}\right)\right]_{0}^{+\infty}
+\int_{0}^{+\infty}\frac{\mathrm{e}^{-pt}}{p}\,\mathrm{d}t\\
&=\frac{1}{p^{2}},
\end{split}
\tag{16.10}
\end{equation}
\begin{equation}
\mathcal{L}\left[t^{2}\right]=\left[t^{2}\left(-\frac{\mathrm{e}^{-pt}}{p}\right)\right]_{0}^{+\infty}
+\frac{2}{p}\int_{0}^{+\infty} t\,\mathrm{e}^{-pt}\,\mathrm{d}t=\frac{2!}{p^{3}}, \tag{16.11}
\end{equation}
and in general (by induction)
\begin{equation}
\mathcal{L}\left[t^{n}\right]=\frac{n!}{p^{\,n+1}},\qquad \operatorname{Re}p>0 . \tag{16.12}
\end{equation}

\subsubsection{Exponential $\mathrm{e}^{at}\leftrightarrow 1/(p-a)$}

\begin{equation}
\mathcal{L}\left[\mathrm{e}^{at}\right]=\int_{0}^{+\infty}\mathrm{e}^{(a-p)t}\,\mathrm{d}t
=\left.\frac{\mathrm{e}^{(a-p)t}}{a-p}\right|_{0}^{+\infty}=\frac{1}{p-a}, \tag{16.13}
\end{equation}
which converges for $\operatorname{Re}(a-p)<0$, i.e.\ $\operatorname{Re}p>\operatorname{Re}a$.

\clearpage
\subsubsection{Sine and cosine (via Euler's formula)}

From $\sin\beta t=\dfrac{\mathrm{e}^{\mathrm{i}\beta t}-\mathrm{e}^{-\mathrm{i}\beta t}}{2\mathrm{i}}$,
and using linearity (Property 1),
\begin{equation}
\mathcal{L}\left[\sin\beta t\right]
=\frac{1}{2\mathrm{i}}\left(\frac{1}{p-\mathrm{i}\beta}-\frac{1}{p+\mathrm{i}\beta}\right)
=\frac{1}{2\mathrm{i}}\cdot\frac{2\mathrm{i}\beta}{p^{2}+\beta^{2}}
=\frac{\beta}{p^{2}+\beta^{2}} . \tag{16.14}
\end{equation}

Similarly, from
$\cos\beta t=\dfrac{\mathrm{e}^{\mathrm{i}\beta t}+\mathrm{e}^{-\mathrm{i}\beta t}}{2}$,
\begin{equation}
\mathcal{L}\left[\cos\beta t\right]=\frac{p}{p^{2}+\beta^{2}} . \tag{16.15}
\end{equation}

\subsubsection{Hyperbolic functions}

From
$\operatorname{sh}\beta t=\dfrac{\mathrm{e}^{\beta t}-\mathrm{e}^{-\beta t}}{2}$ and
$\operatorname{ch}\beta t=\dfrac{\mathrm{e}^{\beta t}+\mathrm{e}^{-\beta t}}{2}$,
\begin{equation}
\mathcal{L}\left[\operatorname{sh}\beta t\right]=\frac{\beta}{p^{2}-\beta^{2}},\qquad
\mathcal{L}\left[\operatorname{ch}\beta t\right]=\frac{p}{p^{2}-\beta^{2}} . \tag{16.16}
\end{equation}

\subsubsection{Table 16.1 --- pairs used in practice}

\begin{center}
\renewcommand{\arraystretch}{1.4}
\begin{tabular}{|c|c|c|c|}
\hline
No. & $f(t)$, $t\ge 0$ & $F(p)=\mathcal{L}[f(t)]$ & ROC \\ \hline
1 & $u(t)$ & $\frac{1}{p}$ & $\operatorname{Re}p>0$ \\ \hline
2 & $t$ & $\frac{1}{p^{2}}$ & $\operatorname{Re}p>0$ \\ \hline
3 & $t^{n}$ & $\frac{n!}{p^{\,n+1}}$ & $\operatorname{Re}p>0$ \\ \hline
4 & $\mathrm{e}^{at}$ & $\frac{1}{p-a}$ & $\operatorname{Re}p>\operatorname{Re}a$ \\ \hline
5 & $\sin\beta t$ & $\frac{\beta}{p^{2}+\beta^{2}}$ & $\operatorname{Re}p>0$ \\ \hline
6 & $\cos\beta t$ & $\frac{p}{p^{2}+\beta^{2}}$ & $\operatorname{Re}p>0$ \\ \hline
7 & $\operatorname{sh}\beta t$ & $\frac{\beta}{p^{2}-\beta^{2}}$ & $\operatorname{Re}p>|\beta|$ \\ \hline
8 & $\operatorname{ch}\beta t$ & $\frac{p}{p^{2}-\beta^{2}}$ & $\operatorname{Re}p>|\beta|$ \\ \hline
\end{tabular}
\end{center}

\noindent\textit{Entry 4 and entries 5--6 are derived with Euler's formula; the rest follow by
direct integration.}

\clearpage
\subsection{Properties 1--4}

Throughout, $\{f(t),F(p)\}$, $\{f_1(t),F_1(p)\}$, $\{f_2(t),F_2(p)\}$ are transform pairs.

\subsubsection{Property 1 --- linearity}

For constants $a,b$,
\begin{equation}
\mathcal{L}\left[a f_1+b f_2\right]=aF_1+bF_2 \tag{16.17}
\end{equation}
and the inverse is linear as well:
\begin{equation}
\mathcal{L}^{-1}\left[aF_1+bF_2\right]=a f_1+b f_2 . \tag{16.18}
\end{equation}

\begin{example}[Partial fractions]
Find $\mathcal{L}^{-1}\left[\frac{p}{(p-1)(p^{2}+4)}\right]$.
\begin{equation*}
\frac{p}{(p-1)(p^{2}+4)}=\frac{A}{p-1}+\frac{Bp+C}{p^{2}+4}
\quad\Longrightarrow\quad A=\frac{1}{5},\qquad Bp+C=-\frac{1}{5}p+\frac{4}{5},
\end{equation*}
so, with Table 16.1 (entries 4, 5, 6),
\begin{equation}
\mathcal{L}^{-1}\left[\frac{p}{(p-1)(p^{2}+4)}\right]
=\frac{1}{5}\mathrm{e}^{t}-\frac{1}{5}\cos 2t+\frac{2}{5}\sin 2t . \tag{16.19}
\end{equation}
\end{example}

\subsubsection{Property 2 --- similarity}

For $a>0$,
\begin{equation}
\mathcal{L}\left[f(at)\right]=\frac{1}{a}F\!\left(\frac{p}{a}\right), \tag{16.20}
\end{equation}
equivalently (replace $a$ by $1/a$),
\begin{equation}
\mathcal{L}\left[\frac{1}{a}f\!\left(\frac{t}{a}\right)\right]=F(ap). \tag{16.21}
\end{equation}

\textit{Proof (sketch).} With $\tau=at$, $\mathrm{d}t=\mathrm{d}\tau/a$,
\begin{equation*}
\mathcal{L}\left[f(at)\right]=\frac{1}{a}\int_{0}^{+\infty} f(\tau)\,\mathrm{e}^{-(p/a)\tau}\,\mathrm{d}\tau
=\frac{1}{a}F\!\left(\frac{p}{a}\right).
\end{equation*}

\clearpage
\subsubsection{Property 3 --- delay (time shift of $f(t)u(t)$)}

\begin{equation}
\mathcal{L}\left[f(t-t_0)\,u(t-t_0)\right]=\mathrm{e}^{-pt_0}F(p), \tag{16.22}
\end{equation}
with inverse
\begin{equation}
\mathcal{L}^{-1}\left[\mathrm{e}^{-pt_0}F(p)\right]=f(t-t_0)\,u(t-t_0)\qquad(t_0>0). \tag{16.23}
\end{equation}

\textbf{Note.} The shift must act on the \emph{whole} function, i.e.\ $f(t-t_0)$ must carry
$u(t-t_0)$; otherwise $f$ is non-zero for $t<t_0$ and the transform is wrong.

\textit{Proof (sketch).} With $\tau=t-t_0$,
\begin{equation*}
\mathcal{L}\left[f(t-t_0)u(t-t_0)\right]
=\int_{t_0}^{+\infty} f(t-t_0)\,\mathrm{e}^{-pt}\,\mathrm{d}t
=\int_{0}^{+\infty} f(\tau)\,\mathrm{e}^{-p(\tau+t_0)}\,\mathrm{d}\tau
=\mathrm{e}^{-pt_0}F(p).
\end{equation*}

\subsubsection{Property 4 --- frequency shift}

\begin{equation}
\mathcal{L}\left[\mathrm{e}^{at}f(t)\right]=F(p-a). \tag{16.24}
\end{equation}

\textit{Proof.} Substitute $p\mapsto p-a$ in the definition:
\begin{equation*}
F(p-a)=\int_{0}^{+\infty} f(t)\,\mathrm{e}^{-(p-a)t}\,\mathrm{d}t
=\int_{0}^{+\infty}\left[\mathrm{e}^{at}f(t)\right]\mathrm{e}^{-pt}\,\mathrm{d}t
=\mathcal{L}\left[\mathrm{e}^{at}f(t)\right].
\end{equation*}

\textbf{Use.} Together with Property 1 and Table 16.1 this is the ``table + partial fractions +
shift'' inversion recipe, e.g.
\begin{equation*}
\mathcal{L}^{-1}\left[\frac{1}{(p-a)^{2}}\right]=t\,\mathrm{e}^{at},\qquad
\mathcal{L}^{-1}\left[\frac{p-a}{(p-a)^{2}+\beta^{2}}\right]=\mathrm{e}^{at}\cos\beta t .
\end{equation*}

\subsubsection{Table 16.2 --- the four properties at a glance}

\begin{center}
\renewcommand{\arraystretch}{1.35}
\begin{tabular}{|c|l|c|}
\hline
Property & Content & No. \\ \hline
1 Linearity & $\mathcal{L}[af_1+bf_2]=a\mathcal{L}[f_1]+b\mathcal{L}[f_2]$; same for $\mathcal{L}^{-1}$ & (16.17)(16.18) \\ \hline
2 Similarity & $\mathcal{L}[f(at)]=\frac{1}{a}F\!\left(\frac{p}{a}\right)$, $a>0$ & (16.20) \\ \hline
3 Delay & $\mathcal{L}[f(t-t_0)u(t-t_0)]=\mathrm{e}^{-pt_0}F(p)$ & (16.22) \\ \hline
4 Frequency shift & $\mathcal{L}[\mathrm{e}^{at}f(t)]=F(p-a)$ & (16.24) \\ \hline
\end{tabular}
\end{center}

\clearpage
\subsection{Homework}

\begin{itemize}
    \item \textbf{A} 9, 10; \textbf{B} 1, 2, 5, 6, 8.
\end{itemize}

\end{document}
